%% =====================================================================
%%  B1-T-02-05 -- what B1 contributes to "Persistence Over Attitude"
%%  -------------------------------------------------------------------
%%  LAYER A: APPEND-ONLY. Written once, then left alone. Material found
%%  only here sits in the `unique` slot and is protected by rule R10.
%% =====================================================================
%% @kind: source
%% @id: B1-T-02-05
%% @book: B1
%% @topic: T-02-05
%% @locator: Tactic 7, pp. 77--88
%% @status: distilled
%% @unique: yes
%% @integrated: yes
%% @updated: 2026-09-19

\dossier{B1}{T-02-05}{\bookshortB1}{Tactic 7, pp. 77--88}

\begin{distilled}
  \item Positive attitude is compared to a gigolo: comfortable and sustaining until hard times arrive, at which point it disappears.
  \item The argument is that you can decide to be persistent and carry the decision out through failure, but you cannot decide to hold an attitude through failure.
  \item Surface attitude --- affirmation routines, self-addressed pep talk --- produces temporary confidence that evaporates on contact with real losses.
  \item Group morale is distinguished from individual stamina: the winning attitude that transforms a team does not carry an individual afterwards.
  \item The chapter's title image contrasts a productive dairy breed with a large impressive-looking one: choose for yield, not for appearance.
\end{distilled}

\begin{unique}
  \item The attitude-versus-persistence argument, including the group-to-individual transfer failure, and the yield-over-appearance selection rule drawn from livestock breeds.
\end{unique}

\begin{terms}
  \item \emph{the Guernsey and the Brahma} --- Sparkman's image for the choice between a productive but unimpressive option and an impressive but unproductive one.
\end{terms}
